\documentclass[aspectratio=169]{beamer}
\usepackage[utf8]{inputenc}
\usepackage[spanish]{babel}
\usepackage{graphicx}
\usepackage{amsmath}
\usepackage{amsfonts}
\usepackage{tikz}
\usetikzlibrary{shapes.geometric, arrows, positioning, calc}
\usepackage{booktabs}

% --- CONFIGURACIÓN DE COLORES INSTITUCIONALES (ESTILO BRAZO ROBÓTICO) ---
\definecolor{USSBlue}{RGB}{0, 32, 91}    % Azul marino principal
\definecolor{USSGold}{RGB}{212, 175, 55} % Dorado de acento
\definecolor{AlertRed}{RGB}{192, 57, 43}  % Rojo alerta
\definecolor{SuccessGreen}{RGB}{39, 174, 96} % Verde ok

\mode<presentation> {
  \usetheme{Madrid}
  \usecolortheme[named=USSBlue]{structure}
  \setbeamercolor{palette primary}{bg=USSBlue,fg=white}
  \setbeamercolor{palette secondary}{bg=USSBlue,fg=white}
  \setbeamercolor{palette tertiary}{bg=USSBlue,fg=white}
  \setbeamercolor{title}{bg=white,fg=USSBlue}
  \setbeamercolor{item}{fg=USSGold}
  \setbeamercolor{block title}{bg=USSBlue,fg=white}
  \setbeamercolor{block body}{bg=USSBlue!5,fg=black}
  \setbeamercolor{block title alert}{bg=AlertRed,fg=white}
  \setbeamercolor{block body alert}{bg=AlertRed!5,fg=black}
  \setbeamertemplate{navigation symbols}{} 
}

% --- DATOS DE PORTADA ---
\title[Proyecto Centinela]{Proyecto Centinela}
\subtitle{Sistema de Alerta Temprana, Monitoreo y Respuesta en Emergencias}
\author[Moisés Amundarain]{Moisés Amundarain}
\institute{Ingeniería Civil Informática}
\date{Julio, 2026}

\begin{document}

% 1. PORTADA
\begin{frame}
  \titlepage
\end{frame}

% 2. DIAGNÓSTICO Y PROBLEMÁTICA
\begin{frame}{¿Por qué nace Centinela? (Lecciones del Temporal)}
  \begin{columns}
    \begin{column}{0.5\textwidth}
      \textbf{El problema real en emergencias:}
      \begin{itemize}
        \item \textbf{Lluvias extremas:} Ríos y quebradas se desbordan de forma rápida.
        \item \textbf{Aislamiento:} Comunidades quedan atrapadas sin poder pedir ayuda.
        \item \textbf{Caminos cortados:} Socavones y aluviones bloquean el paso de vehículos.
      \end{itemize}
    \end{column}
    \begin{column}{0.5\textwidth}
      \textbf{La falla de la tecnología actual:}
      \begin{itemize}
        \item \textbf{Corte de redes:} Las antenas celulares e internet se caen por falta de energía.
        \item \textbf{Corte de agua:} Las plantas potabilizadoras se ensucian y fallan.
        \item \textbf{Desinformación:} Rumores y avisos viejos circulan sin control en redes.
      \end{itemize}
    \end{column}
  \end{columns}
\end{frame}

% 3. ¿CÓMO FUNCIONA CENTINELA?
\begin{frame}{¿Cómo funciona Centinela? (3 Pasos)}
  \centering
  \resizebox{0.92\textwidth}{!}{
    \begin{tikzpicture}[node distance=1.5cm, auto, >=stealth]
      \tikzstyle{bloque1} = [rectangle, draw=USSBlue, fill=USSBlue!10, text width=4.2cm, minimum height=1.2cm, align=center, rounded corners, thick]
      \tikzstyle{bloque2} = [rectangle, draw=USSGold!90!black, fill=USSGold!15, text width=4.2cm, minimum height=1.2cm, align=center, rounded corners, thick]
      \tikzstyle{bloque3} = [rectangle, draw=SuccessGreen, fill=SuccessGreen!10, text width=4.2cm, minimum height=1.2cm, align=center, rounded corners, thick]
      
      \node [bloque1] (c1) {\textbf{1. Sensores y Satélites}\\ \small Miden lluvia, ríos y clima};
      \node [bloque2, right=1.2cm of c1] (c2) {\textbf{2. Análisis Rápido}\\ \small Calcula riesgos de desborde};
      \node [bloque3, right=1.2cm of c2] (c3) {\textbf{3. Alertas y Mapas}\\ \small Avisa a bomberos y vecinos};

      \draw[->, ultra thick, USSBlue] (c1) -- (c2);
      \draw[->, ultra thick, USSGold!90!black] (c2) -- (c3);
    \end{tikzpicture}
  }
  
  \vspace{0.5cm}
  \begin{block}{Principios del Proyecto}
    \begin{itemize}
      \item \textbf{Económico:} Hecho con equipos accesibles de bajo costo.
      \item \textbf{Energía Propia:} Funciona con pequeñas baterías solares.
      \item \textbf{Sin Internet:} Transmite datos aun si se cae la red celular.
    \end{itemize}
  \end{block}
\end{frame}

% 4. INFORMACIÓN SATELITAL
\begin{frame}{Uso de Datos Satelitales}
  \begin{table}
    \centering
  \resizebox{\textwidth}{!}{
    \begin{tabular}{lll}
      \toprule
      \textbf{Tipo de Satélite} & \textbf{Frecuencia de Datos} & \textbf{¿Para qué sirve en Centinela?} \\
      \midrule
      Satélites del Clima (GOES/GPM) & Cada 10 a 30 minutos & Ver la lluvia antes de que caiga en la ciudad. \\
      Satélites de Radar (Sentinel) & Cada 6 días & Ver zonas inundadas a través de las nubes. \\
      Satélites de Mapas (NASA) & Diarios / Horarios & Mapear el avance del agua en el terreno. \\
      \bottomrule
    \end{tabular}
  }
  \end{table}

  \vspace{0.3cm}
  \begin{block}{Ventaja Principal}
    Al cruzar la lluvia del satélite con la nieve en la cordillera, podemos avisar sobre crecidas de ríos con \textbf{varias horas de anticipación}.
  \end{block}
\end{frame}

% 5. ANTES DEL DESASTRE
\begin{frame}{Antes del Desastre: Sensores de Terreno y Radio}
  \begin{columns}
    \begin{column}{0.5\textwidth}
      \textbf{Sensores en Ríos y Quebradas:}
      \begin{itemize}
        \item Miden si el agua sube rápido en puentes o ríos.
        \item Miden la humedad del suelo para prever barro.
        \item Usan carcasas resistentes al agua y sol.
      \end{itemize}
    \end{column}
    \begin{column}{0.5\textwidth}
      \textbf{Transmisión por Radio Mesh:}
      \begin{itemize}
        \item Envía mensajes muy cortos de radio de casa en casa.
        \item No necesita antenas de celular ni cables.
        \item Los datos se empaquetan en códigos pequeños de pocos bytes.
      \end{itemize}
    \end{column}
  \end{columns}

  \vspace{0.4cm}
  \centering
  \resizebox{0.82\textwidth}{!}{
    \begin{tikzpicture}[font=\sffamily\small]
      \draw[fill=USSBlue!20, thick] (0,0) rectangle (2.5,0.8) node[pos=.5] {\textbf{ID del Sensor}};
      \draw[fill=USSGold!30, thick] (2.5,0) rectangle (5.5,0.8) node[pos=.5] {\textbf{Nivel del Agua}};
      \draw[fill=SuccessGreen!20, thick] (5.5,0) rectangle (8.5,0.8) node[pos=.5] {\textbf{Estado Batería}};
      \draw[fill=AlertRed!20, thick] (8.5,0) rectangle (10.5,0.8) node[pos=.5] {\textbf{Alerta}};
    \end{tikzpicture}
  }
\end{frame}

% 6. DESPUÉS DEL DESASTRE: BOTÓN SOS
\begin{frame}{Después del Desastre: Botón SOS desde el Celular}
  \begin{columns}
    \begin{column}{0.5\textwidth}
      \textbf{Página de Emergencia en el Celular:}
      \begin{itemize}
        \item El vecino se conecta al Wi-Fi del sensor local.
        \item El celular obtiene su ubicación por \textbf{GPS satelital} (sin internet ni saldo).
        \item El usuario presiona su estado en 3 opciones:
          \begin{itemize}
            \item \textcolor{AlertRed}{\textbf{[ROJO] Peligro Vital / En el Techo}}
            \item \textcolor{USSGold!90!black}{\textbf{[AMARILLO] Aislado sin Agua}}
            \item \textcolor{SuccessGreen}{\textbf{[VERDE] Aislado pero Bien}}
          \end{itemize}
      \end{itemize}
    \end{column}
    \begin{column}{0.5\textwidth}
      \begin{block}{Mapa de Rescate para Bomberos}
        El mensaje viaja por radio al Puesto de Mando y muestra un mapa con puntos exactos de dónde ir a rescatar primero.
      \end{block}
    \end{column}
  \end{columns}
\end{frame}

% 7. DESPUÉS DEL DESASTRE: PERSONAS ENCONTRADAS
\begin{frame}{Después del Desastre: Personas Encontradas sin Rumores}
  \begin{columns}
    \begin{column}{0.5\textwidth}
      \textbf{El problema de los rumores:}
      \begin{itemize}
        \item En redes sociales circulan afiches de personas perdidas que ya aparecieron.
        \item Esto causa pánico y hace perder tiempo a los equipos de rescate.
      \end{itemize}
    \end{column}
    \begin{column}{0.5\textwidth}
      \textbf{Solución con Código QR:}
      \begin{itemize}
        \item La búsqueda genera un código QR.
        \item Al llegar la persona a un albergue u hospital, se registra en el sistema.
        \item Quien escanee el QR ve de inmediato: \\ \textcolor{SuccessGreen}{\textbf{ENCONTRADO EN ALBERGUE}}.
      \end{itemize}
    \end{column}
  \end{columns}
\end{frame}

% 8. DESPUÉS DEL DESASTRE: AGUA Y RUTAS
\begin{frame}{Después del Desastre: Agua Potable y Rutas Seguras}
  \begin{columns}
    \begin{column}{0.5\textwidth}
      \begin{block}{Sensores en Estanques de Agua}
        \begin{itemize}
          \item Sensores simples en las tapas de los estanques de la comunidad.
          \item Miden cuánta agua queda.
          \item Avisan automáticamente cuando queda menos del 15\% para enviar el camión aljibe.
        \end{itemize}
      \end{block}
    \end{column}
    \begin{column}{0.5\textwidth}
      \begin{block}{Mapa de Rutas Seguras}
        \begin{itemize}
          \item Los vehículos de emergencia reportan caminos cortados o socavones.
          \item El sistema calcula rutas por caminos altos que no se inundan para llevar ayuda.
        \end{itemize}
      \end{block}
    \end{column}
  \end{columns}
\end{frame}

% 9. APOYO A LOS EQUIPOS DE EMERGENCIA
\begin{frame}{Apoyo a los Equipos de Emergencia}
  \centering
  \resizebox{0.82\textwidth}{!}{
    \begin{tikzpicture}[node distance=1.6cm, auto, >=stealth]
      \tikzstyle{centro} = [ellipse, draw=USSBlue, fill=USSBlue!15, text width=3.2cm, minimum height=1.2cm, align=center, thick]
      \tikzstyle{equipo} = [rectangle, draw=black!70, fill=black!5, text width=3cm, minimum height=1cm, align=center, rounded corners]

      \node [centro] (core) {\textbf{Proyecto Centinela}\\ \small Datos y Radio Mesh};
      \node [equipo, above left=0.6cm and 0.8cm of core] (p25) {\textbf{Radios de Bomberos}\\ \small Avisos directos};
      \node [equipo, above right=0.6cm and 0.8cm of core] (cenco) {\textbf{Centrales de Alarma}\\ \small Mapas en tiempo real};
      \node [equipo, below left=0.6cm and 0.8cm of core] (dga) {\textbf{Agua y Caudales}\\ \small Medición de ríos};
      \node [equipo, below right=0.6cm and 0.8cm of core] (senapred) {\textbf{ALERTAS OFICIALES}\\ \small Apoyo a la comunidad};

      \draw[<->, thick, USSBlue] (core) -- (p25);
      \draw[<->, thick, USSBlue] (core) -- (cenco);
      \draw[<->, thick, USSBlue] (core) -- (dga);
      \draw[<->, thick, USSBlue] (core) -- (senapred);
    \end{tikzpicture}
  }
\end{frame}

% 10. CASOS EN LA SERENA Y CONCLUSIÓN
\begin{frame}{Ejemplo Real en La Serena y Próximos Pasos}
  \begin{columns}
    \begin{column}{0.5\textwidth}
      \textbf{Ejemplos en La Serena:}
      \begin{itemize}
        \item \textbf{Pueblo Islón:} Avisar la crecida del río para evacuar de día.
        \item \textbf{Las Compañías:} Recibir llamados SOS por radio sin red de celular.
        \item \textbf{Zona Urbana:} Guiar camiones con agua evitando la Ruta 5 cortada.
      \end{itemize}
    \end{column}
    \begin{column}{0.5\textwidth}
      \textbf{Próximos Pasos:}
      \begin{enumerate}
        \item \textbf{Conversación:} Jjuntar comentarios de amigos bomberos.
        \item \textbf{Simulador:} Probar el sistema primero en computador.
        \item \textbf{Pruebas:} Construir prototipos pequeños para probar en terreno.
      \end{enumerate}
    \end{column}
  \end{columns}

  \vspace{0.5cm}
  \centering
  \textbf{\Large ¡Muchas Gracias!}
\end{frame}

\end{document}
